% Options for packages loaded elsewhere
\PassOptionsToPackage{unicode}{hyperref}
\PassOptionsToPackage{hyphens}{url}
\documentclass[
]{article}
\usepackage{xcolor}
\usepackage[margin=1.5cm]{geometry}
\usepackage{amsmath,amssymb}
\setcounter{secnumdepth}{5}
\usepackage{iftex}
\ifPDFTeX
  \usepackage[T1]{fontenc}
  \usepackage[utf8]{inputenc}
  \usepackage{textcomp} % provide euro and other symbols
\else % if luatex or xetex
  \usepackage{unicode-math} % this also loads fontspec
  \defaultfontfeatures{Scale=MatchLowercase}
  \defaultfontfeatures[\rmfamily]{Ligatures=TeX,Scale=1}
\fi
\usepackage{lmodern}
\ifPDFTeX\else
  % xetex/luatex font selection
\fi
% Use upquote if available, for straight quotes in verbatim environments
\IfFileExists{upquote.sty}{\usepackage{upquote}}{}
\IfFileExists{microtype.sty}{% use microtype if available
  \usepackage[]{microtype}
  \UseMicrotypeSet[protrusion]{basicmath} % disable protrusion for tt fonts
}{}
\makeatletter
\@ifundefined{KOMAClassName}{% if non-KOMA class
  \IfFileExists{parskip.sty}{%
    \usepackage{parskip}
  }{% else
    \setlength{\parindent}{0pt}
    \setlength{\parskip}{6pt plus 2pt minus 1pt}}
}{% if KOMA class
  \KOMAoptions{parskip=half}}
\makeatother
\usepackage{graphicx}
\makeatletter
\newsavebox\pandoc@box
\newcommand*\pandocbounded[1]{% scales image to fit in text height/width
  \sbox\pandoc@box{#1}%
  \Gscale@div\@tempa{\textheight}{\dimexpr\ht\pandoc@box+\dp\pandoc@box\relax}%
  \Gscale@div\@tempb{\linewidth}{\wd\pandoc@box}%
  \ifdim\@tempb\p@<\@tempa\p@\let\@tempa\@tempb\fi% select the smaller of both
  \ifdim\@tempa\p@<\p@\scalebox{\@tempa}{\usebox\pandoc@box}%
  \else\usebox{\pandoc@box}%
  \fi%
}
% Set default figure placement to htbp
\def\fps@figure{htbp}
\makeatother
\setlength{\emergencystretch}{3em} % prevent overfull lines
\providecommand{\tightlist}{%
  \setlength{\itemsep}{0pt}\setlength{\parskip}{0pt}}
\usepackage{booktabs}
\usepackage{float}
\usepackage{placeins}
\usepackage{bookmark}
\IfFileExists{xurl.sty}{\usepackage{xurl}}{} % add URL line breaks if available
\urlstyle{same}
\hypersetup{
  pdftitle={analysis of solar panel data},
  pdfauthor={Hannes Schoss},
  hidelinks,
  pdfcreator={LaTeX via pandoc}}

\title{analysis of solar panel data}
\author{Hannes Schoss}
\date{}

\begin{document}
\maketitle

{
\setcounter{tocdepth}{2}
\tableofcontents
}
i need more detailled time data. I want to evaluate, when the system is
most efficient in a daily thing

There is two ways to maximize this system: One small lever is to cut
down on grid consumption and save a small amount per season roughly this
is about x EUR per year in energy costs.

The much bigger lever is to use up the energy yourself

That is why i want to plot a day from 06:00 to 00:00 in a normal summer
month

x axis: 06:00 -\textgreater{} 00:00

dat{[}dat\$season == ``warm'', {]}

should show a line with the 80\% confidence interval for each of the six
summer months

type: l

\section{solar production over the
day}\label{solar-production-over-the-day}

i do not want the 80\% ci for every month. I want one big ``halo''
80\%ci for the whole warm season, it should be really light and just in
the background, to show the arch

ithere is a battery called battery state of charge. I want a new
variable called first\_full\_battery that marks the time this batttery
is full for the first time in the day

the timespan to this hour should be a grey area in the graphic

`data.frame': 9356 obs. of 15 variables: \$ production : int 6802 6063
1737 4 0 0 0 873 2278 2297 \ldots{} \$ consumption : int 1010 906 756
972 677 346 301 332 879 1333 \ldots{} \$ battery\_charge : int 4022 111
74 0 0 0 0 689 1433 1086 \ldots{} \$ battery\_discharge : int 21 1 10
931 644 313 269 134 48 221 \ldots{} \$ grid\_feedin : int 1819 5047 926
7 2 1 1 20 36 202 \ldots{} \$ grid\_consumption : int 31 0 7 43 36 32 32
35 23 101 \ldots{} \$ battery\_state\_of\_charge: int 66 100 99 85 57 40
29 21 55 87 \ldots{} \$ direct\_consumption : int 987 906 743 4 0 0 0
180 828 1104 \ldots{} \$ time : POSIXct, format: ``2023-03-11 09:00:00''
``2023-03-11 12:00:00'' \ldots{} \$ year : int 2023 2023 2023 2023 2023
2023 2023 2023 2023 2023 \ldots{} \$ month : int 3 3 3 3 3 3 3 3 3 3
\ldots{} \$ day : int 11 11 11 11 11 12 12 12 12 12 \ldots{} \$ hour :
int 9 12 15 18 21 0 3 6 9 12 \ldots{} \$ date : Date, format:
``2023-03-11'' ``2023-03-11'' \ldots{} \$ season : chr ``cold'' ``cold''
``cold'' ``cold'' \ldots{} \textgreater{}

ci\_fun \textless- function(x) \{ n \textless- sum(!is.na(x)) m
\textless- mean(x, na.rm = TRUE) e \textless- qt(0.90, df = n - 1) *
sd(x, na.rm = TRUE) / sqrt(n)

c(mean = m, lower = m - e, upper = m + e) \}

prod\_hour\_month \textless- aggregate( production \textasciitilde{}
month + hour, dat{[}dat\(season == "warm" & dat\)hour \textgreater= 6 \&
dat\$hour \textless= 21, {]}, mean, na.rm = TRUE )

prod\_hour\_warm \textless- aggregate( production \textasciitilde{}
hour, dat{[}dat\(season == "warm" & dat\)hour \textgreater= 6 \&
dat\$hour \textless= 21, {]}, ci\_fun )

prod\_hour\_warm \textless- do.call(data.frame, prod\_hour\_warm)
names(prod\_hour\_warm) \textless- c(``hour'', ``mean'', ``lower'',
``upper'')

plot( NA, xlim = c(6, 21), ylim = range(
c(prod\_hour\_warm\(lower, prod_hour_warm\)upper,
prod\_hour\_month\$production), na.rm = TRUE ), xlab = ``Hour of day'',
ylab = ``Solar production'', xaxt = ``n'', type = ``n'' )

axis( 1, at = c(6, 9, 12, 15, 18, 21), labels = c(``06:00'', ``09:00'',
``12:00'', ``15:00'', ``18:00'', ``21:00'') )

polygon( c(prod\_hour\_warm\(hour, rev(prod_hour_warm\)hour)),
c(prod\_hour\_warm\(lower, rev(prod_hour_warm\)upper)), col = rgb(0, 0,
0, 0.08), border = NA )

lines( prod\_hour\_warm\(hour,
  prod_hour_warm\)mean, lwd = 3, col = rgb(0, 0, 0, 0.35) )

warm\_months \textless- sort(unique(prod\_hour\_month\$month))

for (i in seq\_along(warm\_months)) \{ tmp \textless-
prod\_hour\_month{[}prod\_hour\_month\$month == warm\_months{[}i{]}, {]}

lines( tmp\(hour,
    tmp\)production, lwd = 2, col = i ) \}

legend( ``topright'', legend = month.abb{[}warm\_months{]}, col =
seq\_along(warm\_months), lty = 1, lwd = 2, title = ``Month'' )

\subsection{grid consumption during the
night}\label{grid-consumption-during-the-night}

plot(production \textasciitilde{} hour, dat{[}dat\$season == ``warm'',
{]})

I want another graphic for plotting the grid\_consumption and i want it
to start at 21:00 Uhr and go through the night till the morning at
06:00. I aslo want you to write the version below more concise (e.g.~no
extra dataframe for warm\_dat\ldots.)

warm\_dat \textless- dat{[}dat\$season == ``warm'', {]}

prod\_hour\_month \textless- aggregate( production \textasciitilde{}
month + hour, data = warm\_dat, FUN = function(x) \{ n \textless-
sum(!is.na(x)) m \textless- mean(x, na.rm = TRUE) s \textless- sd(x,
na.rm = TRUE)

\begin{verbatim}
error <- qt(0.90, df = n - 1) * s / sqrt(n)

c(
  mean = m,
  lower = m - error,
  upper = m + error
)
\end{verbatim}

\} )

prod\_hour\_month \textless- do.call( data.frame, prod\_hour\_month )

names(prod\_hour\_month) \textless- c( ``month'', ``hour'', ``mean'',
``lower'', ``upper'' )

prod\_hour\_month \textless- prod\_hour\_month{[}
order(prod\_hour\_month\(month, prod_hour_month\)hour),{]}

plot( NA, xlim = c(6, 24), ylim = range(
c(prod\_hour\_month\(lower, prod_hour_month\)upper), na.rm = TRUE ),
xlab = ``Hour of day'', ylab = ``Production'', xaxt = ``n'', type =
``n'' )

axis( side = 1, at = c(6, 9, 12, 15, 18, 21), labels = c(``06:00'',
``09:00'', ``12:00'', ``15:00'', ``18:00'', ``21:00'') )

warm\_months \textless- sort(unique(prod\_hour\_month\$month))

for (i in seq\_along(warm\_months)) \{ month\_i \textless-
warm\_months{[}i{]} tmp \textless-
prod\_hour\_month{[}prod\_hour\_month\$month == month\_i, {]}

lines( tmp\(hour,
    tmp\)mean, type = ``l'', lwd = 2, col = i )

lines( tmp\(hour,
    tmp\)lower, type = ``l'', lty = 2, col = i )

lines( tmp\(hour,
    tmp\)upper, type = ``l'', lty = 2, col = i ) \}

legend( ``topright'', legend = month.abb{[}warm\_months{]}, col =
seq\_along(warm\_months), lty = 1, lwd = 2, title = ``Month'' )

\begin{enumerate}
\def\labelenumi{\arabic{enumi}.}
\setcounter{enumi}{8}
\tightlist
\item
  Statistical and machine-learning appendix
\end{enumerate}

I would not put this too early. Use it after the main story, as a deeper
analytical layer.

Core questions from your list Which variables explain the largest share
of variance in household grid consumption? Can days be grouped into
typical usage patterns using clustering methods? Can a machine learning
model predict the next day's energy production based on the previous
day's data? Can daily grid consumption be predicted from solar
production, household consumption, and battery behavior? Is there a
statistically significant weekday/weekend difference in solar
production? Is there a statistically significant weekday/weekend
difference in household consumption?

\end{document}
